% GENERATED FILE. Edit canonical lessons, then run npm run generate:books.
% canonical-source-hash: fnv1a64:de104e81d33e6b02
% canonical-lessons: TA-C515-tolilnutpam, TA-C515-menporul, TA-C515-vanporul, TA-C515-matikkanini, TA-C515-minnul

\chapter{Computers}
\label{ch:ta-computers}
\hlchaptermodality{515}

\begin{chapteropening}
\textbf{By the end of this chapter:} \emph{I can talk about computers and what runs on them.}

Complete the last lesson of chapter 515: I can talk about computers and what runs on them.
\end{chapteropening}

\section[\texorpdfstring{\ta{தொழில்நுட்பம்} (toḻilnuṭpam)}{தொழில்நுட்பம் (toḻilnuṭpam)}]{\ta{தொழில்நுட்பம்} (toḻilnuṭpam) — technology}

\label{lesson:TA-C515-tolilnutpam}



\begin{quote}
Before the new one: say the Tamil for a warehouse, then the Tamil for a recipient.
\end{quote}

\subsection*{You'll want to know: \ta{தொழில்நுட்பம்}}
\textbf{\ta{தொழில்நுட்பம்}} — \emph{toḻilnuṭpam} — \textquotedblleft{}technology\textquotedblright{}.

\begin{quote}
\ta{புதிய} \ta{தொழில்நுட்பம்} \ta{விவசாயிகளுக்கு} \ta{உதவுகிறது}.
\end{quote}

\begin{quote}
\emph{putiya toḻilnuṭpam vivacāyikaḷukku utavukiṟatu.}
\end{quote}

\begin{quote}
New technology helps farmers.
\end{quote}

A -\ta{அம்} noun: \ta{தொழில்நுட்பத்தை}, \ta{தொழில்நுட்பத்தில்}.

\begin{cousinweb}
\textbf{\ta{தொழில்நுட்பம்}} (\emph{toḻilnuṭpam}) is a modern compound: the native \ta{தொழில்} (\emph{toḻil}), \textquotedblleft{}work, a craft\textquotedblright{} (DEDR 3524), and \ta{நுட்பம்} (\emph{nuṭpam}), \textquotedblleft{}fineness, subtle skill\textquotedblright{}, which DEDR 3700 lists under \ta{நுண்} (\emph{nuṇ}), \textquotedblleft{}minute, fine\textquotedblright{}.
\end{cousinweb}

\begin{grammarlens}[title={What you've built}]
One more word for describing what you lived and felt.
\end{grammarlens}

\subsection*{Your turn}
\begin{itemize}
\raggedright
  \item \emph{Say it:} \emph{toḻilnuṭpam}
  \item \emph{Say it:} \emph{toḻilnuṭpam}, once more
  \item \emph{Say it:} \emph{putiya toḻilnuṭpam vivacāyikaḷukku utavukiṟatu.} — New technology helps farmers.
  \item \emph{Recall:} say the Tamil for a warehouse, then the Tamil for a recipient, then say \emph{toḻilnuṭpam} again
\end{itemize}

\begin{tcolorbox}[breakable,colback=teal!4,colframe=teal!35!black,title={Before you move on}]
What does \ta{தொழில்நுட்பம்} mean? (\textquotedblleft{}Technology\textquotedblright{}.) Say it once more.
\end{tcolorbox}
\section[\texorpdfstring{\ta{மென்பொருள்} (meṉporuḷ)}{மென்பொருள் (meṉporuḷ)}]{\ta{மென்பொருள்} (meṉporuḷ) — software}

\label{lesson:TA-C515-menporul}



\begin{quote}
Before the new one: say the Tamil for a recipient, then the Tamil for technology.
\end{quote}

\subsection*{You'll want to know: \ta{மென்பொருள்}}
\textbf{\ta{மென்பொருள்}} — \emph{meṉporuḷ} — \textquotedblleft{}software\textquotedblright{}.

\begin{quote}
\ta{நான்} \ta{கணினியில்} \ta{ஒரு} \ta{புதிய} \ta{மென்பொருளைப்} \ta{போட்டேன்}.
\end{quote}

\begin{quote}
\emph{nāṉ kaṇiṉiyil oru putiya meṉporuḷaip pōṭṭēṉ.}
\end{quote}

\begin{quote}
I put a new piece of software on the computer.
\end{quote}

Word for word \textquotedblleft{}soft thing\textquotedblright{}: \ta{மென்} \textquotedblleft{}soft\textquotedblright{} + \ta{பொருள்} (\emph{poruḷ}) \textquotedblleft{}a thing\textquotedblright{}.

\begin{cousinweb}
\textbf{\ta{மென்பொருள்}} (\emph{meṉporuḷ}) is built from native parts to match English \emph{software}: \ta{மென்} (\emph{meṉ}), \textquotedblleft{}soft\textquotedblright{} (DEDR 5078), and \ta{பொருள்} (\emph{poruḷ}), \textquotedblleft{}a thing\textquotedblright{}.
\end{cousinweb}

\begin{grammarlens}[title={What you've built}]
One more word for describing what you lived and felt.
\end{grammarlens}

\subsection*{Your turn}
\begin{itemize}
\raggedright
  \item \emph{Say it:} \emph{meṉporuḷ}
  \item \emph{Say it:} \emph{meṉporuḷ}, once more
  \item \emph{Say it:} \emph{nāṉ kaṇiṉiyil oru putiya meṉporuḷaip pōṭṭēṉ.} — I put a new piece of software on the computer.
  \item \emph{Recall:} say the Tamil for a recipient, then the Tamil for technology, then say \emph{meṉporuḷ} again
\end{itemize}

\begin{tcolorbox}[breakable,colback=teal!4,colframe=teal!35!black,title={Before you move on}]
What does \ta{மென்பொருள்} mean? (\textquotedblleft{}Software\textquotedblright{}.) Say it once more.
\end{tcolorbox}
\section[\texorpdfstring{\ta{வன்பொருள்} (vaṉporuḷ)}{வன்பொருள் (vaṉporuḷ)}]{\ta{வன்பொருள்} (vaṉporuḷ) — hardware}

\label{lesson:TA-C515-vanporul}



\begin{quote}
Before the new one: say the Tamil for technology, then the Tamil for software.
\end{quote}

\subsection*{You'll want to know: \ta{வன்பொருள்}}
\textbf{\ta{வன்பொருள்}} — \emph{vaṉporuḷ} — \textquotedblleft{}hardware\textquotedblright{}.

\begin{quote}
\ta{பிரச்சினை} \ta{மென்பொருளில்} \ta{இல்லை}, \ta{வன்பொருளில்} \ta{இருக்கிறது}.
\end{quote}

\begin{quote}
\emph{piracciṉai meṉporuḷil illai, vaṉporuḷil irukkiṟatu.}
\end{quote}

\begin{quote}
The problem is not in the software, it is in the hardware.
\end{quote}

The pair of \ta{மென்பொருள்} (\emph{meṉporuḷ}): \ta{வன்} \textquotedblleft{}hard\textquotedblright{} + \ta{பொருள்} \textquotedblleft{}a thing\textquotedblright{}.

\begin{cousinweb}
\textbf{\ta{வன்பொருள்}} (\emph{vaṉporuḷ}) is built from native parts to match English \emph{hardware}: \ta{வன்} (\emph{vaṉ}), \textquotedblleft{}hard, strong\textquotedblright{} (DEDR 5276), and \ta{பொருள்} (\emph{poruḷ}), \textquotedblleft{}a thing\textquotedblright{}.
\end{cousinweb}

\begin{grammarlens}[title={What you've built}]
One more word for describing what you lived and felt.
\end{grammarlens}

\subsection*{Your turn}
\begin{itemize}
\raggedright
  \item \emph{Say it:} \emph{vaṉporuḷ}
  \item \emph{Say it:} \emph{vaṉporuḷ}, once more
  \item \emph{Say it:} \emph{piracciṉai meṉporuḷil illai, vaṉporuḷil irukkiṟatu.} — The problem is not in the software, it is in the hardware.
  \item \emph{Recall:} say the Tamil for technology, then the Tamil for software, then say \emph{vaṉporuḷ} again
\end{itemize}

\begin{tcolorbox}[breakable,colback=teal!4,colframe=teal!35!black,title={Before you move on}]
What does \ta{வன்பொருள்} mean? (\textquotedblleft{}Hardware\textquotedblright{}.) Say it once more.
\end{tcolorbox}
\section[\texorpdfstring{\ta{மடிக்கணினி} (maṭikkaṇiṉi)}{மடிக்கணினி (maṭikkaṇiṉi)}]{\ta{மடிக்கணினி} (maṭikkaṇiṉi) — a laptop}

\label{lesson:TA-C515-matikkanini}



\begin{quote}
Before the new one: say the Tamil for software, then the Tamil for hardware.
\end{quote}

\subsection*{You'll want to know: \ta{மடிக்கணினி}}
\textbf{\ta{மடிக்கணினி}} — \emph{maṭikkaṇiṉi} — \textquotedblleft{}a laptop\textquotedblright{}.

\begin{quote}
\ta{அவன்} \ta{பயணத்தில்} \ta{மடிக்கணினியைக்} \ta{கொண்டு} \ta{போனான்}.
\end{quote}

\begin{quote}
\emph{avaṉ payaṇattil maṭikkaṇiṉiyaik koṇṭu pōṉāṉ.}
\end{quote}

\begin{quote}
He took his laptop on the journey.
\end{quote}

Word for word \textquotedblleft{}lap computer\textquotedblright{}: \ta{மடி} \textquotedblleft{}a lap\textquotedblright{} + \ta{கணினி} (\emph{kaṇiṉi}).

\begin{cousinweb}
\textbf{\ta{மடிக்கணினி}} (\emph{maṭikkaṇiṉi}) is built from native Tamil parts to match English \emph{laptop}: \ta{மடி} (\emph{maṭi}), \textquotedblleft{}a lap\textquotedblright{}, and \ta{கணினி} (\emph{kaṇiṉi}), \textquotedblleft{}a computer\textquotedblright{}; \ta{மடி} is usually traced to the native verb \ta{மடி}, \textquotedblleft{}to fold, to be lapped in\textquotedblright{} (DEDR 4645).
\end{cousinweb}

\begin{grammarlens}[title={What you've built}]
One more word for describing what you lived and felt.
\end{grammarlens}

\subsection*{Your turn}
\begin{itemize}
\raggedright
  \item \emph{Say it:} \emph{maṭikkaṇiṉi}
  \item \emph{Say it:} \emph{maṭikkaṇiṉi}, once more
  \item \emph{Say it:} \emph{avaṉ payaṇattil maṭikkaṇiṉiyaik koṇṭu pōṉāṉ.} — He took his laptop on the journey.
  \item \emph{Recall:} say the Tamil for software, then the Tamil for hardware, then say \emph{maṭikkaṇiṉi} again
\end{itemize}

\begin{tcolorbox}[breakable,colback=teal!4,colframe=teal!35!black,title={Before you move on}]
What does \ta{மடிக்கணினி} mean? (\textquotedblleft{}A laptop\textquotedblright{}.) Say it once more.
\end{tcolorbox}
\section[\texorpdfstring{\ta{மின்னூல்} (miṉṉūl)}{மின்னூல் (miṉṉūl)}]{\ta{மின்னூல்} (miṉṉūl) — an e-book}

\label{lesson:TA-C515-minnul}



\begin{quote}
Before the new one: say the Tamil for hardware, then the Tamil for a laptop.
\end{quote}

\subsection*{You'll want to know: \ta{மின்னூல்}}
\textbf{\ta{மின்னூல்}} — \emph{miṉṉūl} — \textquotedblleft{}an e-book\textquotedblright{}.

\begin{quote}
\ta{நான்} \ta{இப்போது} \ta{மின்னூல்களைக்} \ta{கைப்பேசியில்} \ta{படிக்கிறேன்}.
\end{quote}

\begin{quote}
\emph{nāṉ ippōtu miṉṉūlkaḷaik kaippēciyil paṭikkiṟēṉ.}
\end{quote}

\begin{quote}
These days I read e-books on my phone.
\end{quote}

Plural: \ta{மின்னூல்கள்}, \emph{miṉṉūlkaḷ}.

\ta{மின்} \textquotedblleft{}electric\textquotedblright{} + \ta{நூல்} \textquotedblleft{}a book\textquotedblright{} (a literary word for \ta{புத்தகம்}).

\begin{cousinweb}
\textbf{\ta{மின்னூல்}} (\emph{miṉṉūl}) is built from native parts: \ta{மின்} (\emph{miṉ}), \textquotedblleft{}a flash, lightning\textquotedblright{}, now \textquotedblleft{}electric\textquotedblright{} (DEDR 4876), and \ta{நூல்} (\emph{nūl}), \textquotedblleft{}a book\textquotedblright{}; after \ta{ன்} the \ta{ந} of \ta{நூல்} is written \ta{ன}: \ta{மின்} + \ta{நூல்} = \ta{மின்னூல்}.
\end{cousinweb}

\begin{grammarlens}[title={What you've built}]
One more word for describing what you lived and felt.
\end{grammarlens}

\subsection*{Your turn}
\begin{itemize}
\raggedright
  \item \emph{Say it:} \emph{miṉṉūl}
  \item \emph{Say it:} \emph{miṉṉūl}, once more
  \item \emph{Say it:} \emph{nāṉ ippōtu miṉṉūlkaḷaik kaippēciyil paṭikkiṟēṉ.} — These days I read e-books on my phone.
  \item \emph{Recall:} say the Tamil for hardware, then the Tamil for a laptop, then say \emph{miṉṉūl} again
\end{itemize}

\begin{tcolorbox}[breakable,colback=teal!4,colframe=teal!35!black,title={Before you move on}]
What does \ta{மின்னூல்} mean? (\textquotedblleft{}An e-book\textquotedblright{}.) Say it once more.
\end{tcolorbox}
